\section{Theoretical Setup}\label{sec:setup}

\subsection{Problem Setting}

Let \(\log\) denote the natural logarithm.  Let \(\log^*\) denote
base-two iterated logarithm, with \(\log^*(t)\) the least
\(j\geq0\) for which the \(j\)-fold iterated base-two logarithm of
\(t\) is at most \(1\).  For a domain \(X\) and a nonempty finite
binary concept class \(C\subseteq\{0,1\}^X\), set
\[
d(C)=\operatorname{VC}(C),\qquad
\ell(C)=\operatorname{LD}(C),\qquad
s(C)=1+\log^*(1+\ell(C)),\qquad
q(C)=d(C)+s(C).
\]
Thus \(q(C)\geq1\).  Finiteness gives the counting bound
\[
\ell(C)\leq\log_2|C|.
\]

Write \(Z_X=X\times\{0,1\}\) and
\(\mathcal H_X=\{0,1\}^X\).  A distribution \(Q\) on \(Z_X\) is
realizable by \(C\) if some \(c\in C\) satisfies
\(Q\{(x,y):y=c(x)\}=1\).  For \(h\in\mathcal H_X\), define
\[
R_Q(h)=\Pr_{(x,y)\sim Q}[h(x)\neq y].
\]
All randomized algorithms are kernels on output spaces in which the
singleton hypotheses and output events used below are measurable.
An equality \(G(S)=h\) is equality as functions on \(X\), and its
probability includes both the iid sample and all internal randomness.

Two ordered samples of equal size are replacement-adjacent if they
differ in at most one record.  A kernel
\(A:Z_X^n\rightsquigarrow\mathcal H_X\) is
\((\varepsilon,\delta)\)-differentially private if, for every such
\(S,S'\) and every measurable output event \(E\),
\[
\Pr[A(S)\in E]
\leq e^\varepsilon\Pr[A(S')\in E]+\delta,
\]
and the same inequality holds after interchanging \(S,S'\).  It is
distribution-free realizable \((\alpha,\beta)\)-PAC for \(C\) if,
for every \(Q\) realizable by \(C\),
\[
\Pr_{S\sim Q^n,\,h\sim A(S)}[R_Q(h)>\alpha]\leq\beta.
\]
There is no properness, output-representation, or computational
restriction.  Fix
\[
\alpha_0=\frac18,\qquad \beta_0=\frac18,
\]
and regard \(\varepsilon_0\in(0,1)\) as an arbitrary fixed constant.

\subsection{Assumption}

\begin{assumption}[Polynomial global-stability profile]
\label{assump:polynomial-global-stability}
There is a universal integer \(a\geq1\) such that, for every domain
\(X\) and every nonempty finite \(C\subseteq\{0,1\}^X\), there are an
integer
\[
1\leq m_C\leq q(C)^a
\]
and a single randomized nonprivate producer
\[
G_C:Z_X^{m_C}\rightsquigarrow\mathcal H_X,
\]
chosen independently of the realizable distribution, with the
following property.  For every distribution \(Q\) realizable by
\(C\), there is \(h_{C,Q}\in\mathcal H_X\) such that
\[
R_Q(h_{C,Q})\leq\frac{\alpha_0}{2},
\qquad
\Pr_{S\sim Q^{m_C},\,G_C}[G_C(S)=h_{C,Q}]
\geq q(C)^{-a}.
\tag{GS}
\]
The atom may depend on \(Q\), but \(G_C,m_C\), and \(a\) may not.
Both the producer and the atom may be improper and computationally
unbounded.
\end{assumption}

This atom condition is the sole nonstandard primitive hypothesis.
Private learnability, the stable-selection conversion, its
quantitative specialization, and all sequence asymptotics are
conclusions to be proved.

\subsection{Formalized Goal}

For a sequence of nonempty finite classes
\(\{C_\kappa\}_{\kappa\in\mathbb N}\), write
\(L_\kappa=\log|C_\kappa|\) and use the corresponding subscripts on
\(d,\ell,s,q\).  The target is the following conditional statement.
For every fixed \(\varepsilon_0\in(0,1)\), every sequence satisfying
\[
|C_\kappa|\longrightarrow\infty
\]
and
\[
\forall p\in\mathbb N\ \exists\kappa_0(p)\
\forall\kappa\geq\kappa_0(p):
\qquad L_\kappa>d_\kappa^p
\tag{SP}
\]
has, at every index, an arbitrary-output, computationally
unrestricted, distribution-free realizable
\((\alpha_0,\beta_0)\)-PAC learner that is
\((\varepsilon_0,\delta_\kappa)\)-differentially private and uses
exactly
\[
N_\kappa=N_a(q_\kappa,L_\kappa,\varepsilon_0)
\]
records, where the functions \(N_a\) and \(\delta_a\) are defined in
Section~\ref{sec:prelim}.  The required conclusions are
\[
N_\kappa=L_\kappa^{o(1)}=o(L_\kappa),
\qquad
\delta_\kappa N_\kappa^\rho\longrightarrow0
\quad\text{for every fixed }\rho>0.
\]
The same learner must remain eligible whenever an allowed privacy
schedule is pointwise no smaller than \(\delta_\kappa\), including
the standard allowance
\(c/[N_\kappa^2\log(eN_\kappa)]\) eventually for every fixed
\(c>0\).  Consequently no eventual positive
\(\Omega(L_\kappa)\) lower bound can hold on this unrestricted learner
interface.  The conclusion remains conditional on
Assumption~\ref{assump:polynomial-global-stability}; the assumption
itself is not asserted or proved.
